% Source: Wald GR, physical PDF p. 237 (printed p. 224), Figure 9.2.
% Coverage: Conjugate points along a timelike geodesic.
% Figures/tables/footnotes: Figure 9.2.
% Status: source-reviewed and compile-clean.
% Uncertainties: none.

\begingroup
\renewcommand{\thefigure}{9.2}
\begin{figure}[tb]
  \centering
  \begin{tikzpicture}[x=1cm,y=1cm,>=Stealth,thick]
    \coordinate (p) at (0,0);
    \coordinate (q) at (1.2,4.5);
    \draw (p) .. controls (-.1,1.3) and (.6,3.2) .. (q);
    \draw (p) .. controls (.42,1.35) and (1.35,3.2) .. (q);
    \fill (p) circle (2pt) node[below right] {\(p\)};
    \fill (q) circle (2pt) node[above right] {\(q\)};
    \draw[->] (1.55,2.62) -- (.73,2.62);
    \node[right] at (1.55,2.62) {\(\gamma\)};
    \draw[->] (.61,1.95) -- (.18,1.95);
    \node[left] at (.13,1.95) {\(\eta^a\)};
    % Reserve trim clearance above the artwork when this float lands at page top.
    \path[draw=none] (0,5.5) circle[radius=.1pt];
  \end{tikzpicture}
  \caption{A spacetime diagram illustrating the notion of conjugate points along the
  geodesic \(\gamma\).}
  \label{fig:9.2}
\end{figure}
\endgroup
